% actividad-biblioteca-sesion01.tex
% Actividad de investigación en biblioteca — Sesión 1 (Unidad I, tema 1.1)
% Laboratorio: Gestión de datos, comportamiento y perfiles del consumidor digital
% Lic. en Mercadotecnia Digital · ESCA Unidad Santo Tomás · IPN
\documentclass[11pt,letterpaper]{article}

\usepackage[utf8]{inputenc}
\usepackage[spanish,es-noshorthands]{babel}
\usepackage[T1]{fontenc}
\usepackage{geometry}
\geometry{top=2.5cm, bottom=2.5cm, left=2.5cm, right=2.5cm}
\usepackage{csquotes}
\usepackage{enumitem}
\usepackage{booktabs}
\usepackage{array}
\usepackage{amssymb}
\usepackage{xcolor}
\usepackage{hyperref}
\usepackage{fancyhdr}
\usepackage{tcolorbox}
\tcbuselibrary{skins,breakable}

\definecolor{guinda}{HTML}{6F1D46}

\hypersetup{
  colorlinks=true, linkcolor=guinda, urlcolor=guinda,
  pdftitle={Actividad de investigación en biblioteca -- Sesión 1},
  pdfauthor={Dr. Aboud Barsekh-Onji}
}

\setlength{\headheight}{14pt}
\pagestyle{fancy}
\fancyhf{}
\lhead{\small Laboratorio: Gestión de Datos, Comportamiento y Perfiles del Consumidor Digital}
\lfoot{\small IPN --- ESCA Unidad Santo Tomás}
\rfoot{\small Página \thepage}
\setlength{\parindent}{0pt}
\setlength{\parskip}{0.5\baselineskip}

\newtcolorbox{caja}[1][]{colback=guinda!6, colframe=guinda, boxrule=0.6pt, arc=1mm,
  left=6pt, right=6pt, top=5pt, bottom=5pt, breakable, #1}

\begin{document}
\thispagestyle{fancy}

\begin{center}
  {\LARGE \textbf{Investigación en biblioteca:}}\\[3pt]
  {\LARGE \textbf{del relato editorial a la evidencia verificable}}\\[6pt]
  {\large Actividad de la Sesión 1 --- Unidad I, tema 1.1 \emph{Comportamiento del consumidor}}\\[4pt]
  {\small Lic. en Mercadotecnia Digital $\cdot$ ESCA Unidad Santo Tomás $\cdot$ IPN}\\[2pt]
  {\small Dr. Aboud Barsekh-Onji}
\end{center}

\begin{tabular}{@{}lp{10.5cm}@{}}
  \textbf{Nombre(s):} & \rule{9cm}{0.4pt} \\[4pt]
  \textbf{Equipo / grupo:} & \rule{9cm}{0.4pt} \\[4pt]
  \textbf{Modalidad:} & Equipos de 3 a 4 personas \\
  \textbf{Duración estimada:} & 1 visita a biblioteca (2 h) + 3 h de análisis y redacción \\
  \textbf{Profesor:} & Dr. Aboud Barsekh-Onji \\
  \textbf{Fecha límite de entrega:} & \textbf{Viernes 2 de octubre} a más tardar, en la carpeta de evidencias (portafolio) \\
  \textbf{Producto:} & Reporte de investigación de 4 a 6 páginas + ficha de fuentes \\
\end{tabular}

\vspace{0.2cm}

\section*{1. Propósito}

En la sesión 1 distinguimos dos relatos sobre el consumidor digital: una
\textbf{curva medible} (estadística pública) y una \textbf{narrativa editorial}
(la serie \textit{Marketing X.0}). También vimos que \enquote{inteligente} no
equivale a \enquote{sostenible}, y que las cuatro ventajas del análisis del
consumidor tienen niveles de evidencia desiguales. Esta actividad les pide
salir de las diapositivas y \textbf{contrastar esas afirmaciones contra las
fuentes originales} disponibles en la biblioteca.

\begin{caja}
\textbf{Resultados de aprendizaje.} Al terminar, cada equipo podrá:
\begin{enumerate}[leftmargin=1.4em, itemsep=2pt]
  \item Localizar y distinguir fuentes primarias, académicas, de industria y estadísticas.
  \item Evaluar la solidez de una afirmación sobre el consumidor digital según su fuente.
  \item Contrastar la periodización editorial (X.0) con evidencia empírica verificable.
  \item Documentar fuentes con rigor: autor, año, edición, página y fecha de consulta.
\end{enumerate}
\end{caja}

\section*{2. Ejes de investigación}

Cada equipo trabaja \textbf{un eje} (se asigna en clase; no se repiten ejes dentro del grupo).

\begin{center}
\small
\begin{tabular}{@{}>{\raggedright}p{1.1cm}>{\raggedright}p{4.6cm}>{\raggedright\arraybackslash}p{8.4cm}@{}}
\toprule
\textbf{Eje} & \textbf{Tema de la sesión} & \textbf{Pregunta rectora} \\
\midrule
A & 1.1.1 Evolución del consumidor & ¿Qué evidencia sostiene que el consumidor \enquote{cambió de versión} (3.0, 4.0, 5.0, 6.0)? ¿Qué cambió de forma verificable? \\
B & 1.1.2 Consumo inteligente y sostenible & ¿Qué dice la literatura sobre la brecha actitud-conducta en consumo sostenible y qué se sabe de ella en México? \\
C & 1.1.3 Ventajas del análisis & ¿Qué respaldo empírico tiene cada ventaja del análisis del consumidor (segmentar, personalizar, reducir fricción, anticipar tendencias)? \\
D & Evidencia de México & ¿Qué muestran las series estadísticas de INEGI y de la industria sobre el consumidor digital mexicano, y con qué límites metodológicos? \\
\bottomrule
\end{tabular}
\end{center}

\section*{3. Desarrollo}

\subsection*{Fase 1 --- Planeación (antes de ir a la biblioteca, 30 min)}
\begin{enumerate}[leftmargin=1.4em, itemsep=3pt]
  \item Reformulen la pregunta rectora de su eje con sus propias palabras.
  \item Escriban \textbf{tres palabras clave} en español y tres en inglés, y sus sinónimos.
  \item Propongan una \textbf{hipótesis inicial} (una frase) que la investigación pondrá a prueba.
\end{enumerate}

\subsection*{Fase 2 --- Búsqueda en biblioteca (2 h)}
\begin{enumerate}[leftmargin=1.4em, itemsep=3pt]
  \item Consulten el catálogo de las bibliotecas del IPN (\url{https://www.busqueda.dirbibliotecas.ipn.mx}) y sus bases de datos o repositorios institucionales.
  \item Localicen \textbf{al menos 5 fuentes}, con esta composición mínima:
\end{enumerate}

\begin{center}
\small
\begin{tabular}{@{}>{\raggedright}p{6.2cm}>{\raggedright\arraybackslash}p{7.6cm}@{}}
\toprule
\textbf{Tipo de fuente} & \textbf{Mínimo requerido} \\
\midrule
Libro de texto o de autor de referencia (impreso o digital) & 1 \\
Artículo académico arbitrado (revista con revisión por pares) & 2 \\
Fuente estadística oficial (por ejemplo, INEGI) o de industria con metodología pública & 1 \\
Fuente institucional o de política pública (por ejemplo, PROFECO, ONU) & 1 \\
\bottomrule
\end{tabular}
\end{center}

\begin{caja}
\textbf{Punto de partida (no es lista cerrada).} Estas obras se citaron en la
sesión 1 y pueden servirles para arrancar: Kotler, Kartajaya y Setiawan
(\textit{Marketing 5.0} y \textit{6.0}); Solomon (\textit{Comportamiento del
consumidor}); Sheth y Jain (\textit{Digital consumer behavior}); Smith (1956);
Bass (1969); Tversky y Kahneman (1974); INEGI (ENDUTIH); AMVO; PROFECO.
\textbf{Su reto es encontrar fuentes que \emph{no} se citaron en clase.}
\end{caja}

\subsection*{Fase 3 --- Lectura crítica y ficha de fuentes (1.5 h)}
Por cada fuente, llenen una ficha con los campos de la tabla. El campo
\textbf{nivel de evidencia} usa la escala de la sesión: \emph{alta},
\emph{media}, \emph{baja-media}.

\begin{center}
\small
\begin{tabular}{@{}>{\raggedright}p{4.3cm}>{\raggedright\arraybackslash}p{9.5cm}@{}}
\toprule
\textbf{Campo de la ficha} & \textbf{Qué deben registrar} \\
\midrule
Referencia completa (APA 7) & Autor(es), año, título, editorial o revista, volumen, páginas, DOI o URL \\
Ubicación & Signatura topográfica, base de datos o repositorio donde la localizaron \\
Tipo de fuente & Libro, artículo arbitrado, estadística oficial, institucional, industria \\
Dato o idea central & Una afirmación concreta, con \textbf{página} o tabla de origen \\
Método y datos & ¿Qué muestra, periodo y método usa? ¿Es replicable? \\
Nivel de evidencia & Alta / media / baja-media, con \textbf{una frase} de justificación \\
Límites y sesgos & Qué no permite afirmar la fuente; posible interés del autor o editor \\
\bottomrule
\end{tabular}
\end{center}

\subsection*{Fase 4 --- Reporte (1.5 h)}
El reporte de 4 a 6 páginas incluye las siguientes secciones:
\begin{enumerate}[leftmargin=1.4em, itemsep=3pt]
  \item \textbf{Pregunta e hipótesis inicial} (media página).
  \item \textbf{Hallazgos:} qué encontraron, organizado por argumento y no por fuente.
  \item \textbf{Contraste con la sesión 1:} ¿confirman, matizan o contradicen lo visto en clase? Cítenlo con precisión.
  \item \textbf{Cuadro comparativo} de al menos tres fuentes según su nivel de evidencia.
  \item \textbf{Conclusión y hueco de investigación:} qué sigue sin respuesta o carece de dato para México.
  \item \textbf{Referencias} en formato APA 7 y \textbf{anexo} con las fichas de fuentes.
\end{enumerate}

\begin{caja}
\textbf{Reglas de rigor (no negociables).}
\begin{itemize}[leftmargin=1.4em, itemsep=2pt]
  \item Toda cifra lleva fuente, año y \textbf{fecha de consulta}. Sin fuente, no se usa.
  \item Si no encuentran un dato, escriban \enquote{no localizado}: \textbf{no lo inventen ni lo estimen de memoria}.
  \item No se acepta una referencia que no hayan abierto. Cada ficha exige página o tabla de origen.
  \item Si usan herramientas de IA para orientar su búsqueda, verifiquen cada referencia en el catálogo o en la fuente: las referencias inexistentes o alteradas se califican como cero en ese criterio.
\end{itemize}
\end{caja}

\section*{4. Rúbrica de evaluación (100 puntos)}

\begin{center}
\small
\begin{tabular}{@{}>{\raggedright}p{5.2cm}>{\raggedright}p{6.6cm}r@{}}
\toprule
\textbf{Criterio} & \textbf{Evidencia esperada} & \textbf{Pts.} \\
\midrule
Calidad y variedad de fuentes & Cumple la composición mínima; fuentes verificables y pertinentes al eje & 20 \\
Fichas de fuentes & Campos completos, con página y nivel de evidencia justificado & 20 \\
Análisis crítico & Distingue dato, interpretación y opinión; identifica límites y sesgos & 25 \\
Contraste con la sesión 1 & Vincula hallazgos con conceptos de clase con precisión & 15 \\
Redacción y citación APA 7 & Claridad, coherencia, referencias correctas y consistentes & 10 \\
Honestidad intelectual & Sin datos sin fuente ni referencias no verificadas; hueco de investigación explícito & 10 \\
\midrule
\textbf{Total} & & \textbf{100} \\
\bottomrule
\end{tabular}
\end{center}

\section*{5. Lista de cotejo antes de entregar}

\begin{itemize}[leftmargin=1.6em, label=$\square$, itemsep=2pt]
  \item Mi equipo trabajó el eje asignado y su pregunta rectora aparece en el reporte.
  \item Hay al menos 5 fuentes con la composición mínima y al menos una que no se citó en clase.
  \item Cada fuente tiene ficha completa, con página o tabla de origen.
  \item Cada cifra del reporte tiene fuente, año y fecha de consulta.
  \item El reporte contrasta explícitamente con la sesión 1 e incluye un hueco de investigación.
  \item Las referencias están en APA 7 y el archivo se subió a la carpeta del equipo en el portafolio.
\end{itemize}

\vspace{0.3cm}
\begin{center}
\small\textit{Pregunta para cerrar su reporte: ¿qué le preguntarían a esta fuente antes de creerle?}
\end{center}

\end{document}
